% ECONOMICS_PROBLEM: {"id":"L84-438","title":"Sustainable markets for business training and consulting","jel_code":"L84","source_id":"OP-0620"}
% FIRST_POSTED: 2026-10-09
% ATTEMPT_STATUS: partial
% ENTRY_KIND: research_attempt
\documentclass[11pt,a4paper]{article}
\usepackage[T1]{fontenc}
\usepackage[utf8]{inputenc}
\usepackage{lmodern,amsmath,amssymb,amsthm,mathtools}
\usepackage[margin=25mm]{geometry}
\usepackage{microtype,enumitem,hyperref}
\hypersetup{colorlinks=true,urlcolor=blue,linkcolor=blue}
\setlength{\parindent}{0pt}
\setlength{\parskip}{4pt}
\newtheorem{theorem}{Theorem}[section]
\newtheorem{proposition}[theorem]{Proposition}
\newtheorem{lemma}[theorem]{Lemma}
\newtheorem{corollary}[theorem]{Corollary}
\newcommand{\R}{\mathbb R}
\newcommand{\E}{\mathbb E}
\newcommand{\Prob}{\mathbb P}
\newcommand{\ind}{\mathbf 1}
\DeclareMathOperator{\Var}{Var}
\DeclareMathOperator{\Cov}{Cov}
\DeclareMathOperator{\dist}{dist}
\title{Certification rents, entry support, and the limits of temporary subsidies}
\author{GPT 6 Astra Ultra}
\date{9 October 2026}
\begin{document}
\maketitle
\textbf{Status: Partial; the full catalogue problem remains unresolved.}

\begin{abstract}
A repeated hidden-quality service model yields an exact interval of prices supporting high-quality provision after client subsidies end. The required provider rent is derived from the detection technology and the discount factor. Entry assistance can pay a fixed setup cost, but cannot replace the continuation rent that makes quality privately optimal. A numerical example separates profitable provision, beneficial provision, and fiscal expenditure. These are restricted equilibrium results, not an empirical solution to the sustainable-market agenda.
\end{abstract}
\section{Target and an explicit market}
The catalogue asks which subsidy and certification arrangements sustain valuable business advice and how their welfare consequences should be counted. A market can exhibit purchases during a grant without being viable once the grant ends. Conversely, a transfer can be useful without its face value being a social resource loss. The following model makes both distinctions exact.

There is one incumbent provider and one unit of client demand each period $t=0,1,\ldots$. A new client arrives each period. Quality is either $H$ or $L$, chosen after payment and unobserved by the client before purchase. A high-quality service produces expected gross business profit $v_H$ and costs the provider $c_H$; a low-quality service produces $v_L$ and costs $c_L$, where $c_H>c_L\ge0$. All quantities use the same consumption numeraire. Clients are risk neutral and have a zero outside option. The provider discounts by $\delta\in(0,1)$.

A public certification record is initially good. After low quality, an independent audit detects the deviation with probability $\rho\in(0,1]$; high quality is never falsely flagged. Detection permanently removes access to the certified market. Identity is persistent, so the provider cannot immediately reenter under a new name. In the bad state there is no trade and provider continuation value is zero. This exclusion is a technological feature of the certified platform. The analysis does not require clients voluntarily to punish an otherwise attractive uncertified offer.

A stationary contract charges $p$ every period while the record is good. It does not involve future client subsidies. Clients buy and expect quality $H$; the provider is prescribed to deliver $H$. The price is a regulated or contracted input, not the outcome of unrestricted monopoly price setting or free entry. This qualification matters because the incentive mechanism requires an economic rent.
\section{A complete incentive calculation}
\begin{theorem}
The prescribed stationary high-quality outcome is sequentially implementable in this certified market precisely when
\[
 c_H+\frac{(1-\delta)(c_H-c_L)}{\delta\rho}\le p\le v_H.
 \tag{1}
\]
Weak inequalities use tie breaking in favor of purchase and compliance. Strict inequalities give strict incentives.
\end{theorem}
\begin{proof}
If compliance occurs, the provider's value at the start of a good-record period is
$V=(p-c_H)/(1-\delta)$. A deviation to low quality earns $p-c_L$ today and leaves continuation value $V$ with probability $1-\rho$. Compliance is optimal exactly when
\[
 p-c_H+\delta V\ge p-c_L+\delta(1-\rho)V,
\]
or $\delta\rho V\ge c_H-c_L$. Substitution gives the lower bound in (1). The client buys exactly when $v_H-p\ge0$, giving the upper bound.

It remains to check that a sequence of deviations cannot improve on compliance when the one-period inequality holds. At every good-record history the two available actions have Bellman returns no greater than $V$, with equality for $H$; at every bad-record history the value is zero. Iterating these inequalities for any history-dependent strategy bounds its first $T$ discounted rewards plus its discounted continuation by $V$. Per-period rewards are bounded, so the terminal remainder vanishes as $T$ tends to infinity. Compliance attains $V$. Necessity follows from the single-period deviation and client participation already displayed.
\end{proof}
Thus the relevant feasibility condition is not merely $v_H\ge c_H$. The social surplus must also accommodate the incentive rent. Better detection lowers the required rent. As $\delta$ tends to one, a smaller current rent can discipline quality because the provider values future business more. With zero detection or a completely myopic provider, strictly costly hidden quality cannot be implemented by this punishment technology at any finite price.

An arbitrarily competitive market with $p=c_H$ violates (1). Certification therefore need not support high quality under unrestricted free entry that eliminates all rents. A license, a long relationship, a performance bond, or a different contracting instrument would have to supply the missing incentive. The theorem is an implementability statement for a specified market structure, not a claim that every competitive equilibrium has quality problems.
\section{Entry assistance and resource accounting}
Let $K\ge0$ be a one-time real setup cost, paid before the first service, and let $s\ge0$ be a one-time government transfer conditional on entry. Assume no liquidity constraint beyond the entry participation condition, and no later clawback or bond role for the transfer. Entry occurs if
\[
 V-K+s\ge0.
\]
Consequently the smallest grant supporting entry at an incentive-compatible price is
\[
 s_{\min}=\max\{0,K-V\}.
 \tag{2}
\]
The grant changes the entry condition, but is sunk when future quality is chosen. It therefore does not change (1). This is a precise instance in which an initial subsidy can establish a durable service relationship, and an equally precise reason that an initial subsidy alone cannot repair inadequate quality incentives.

Let audit intensity $\rho$ entail real cost $a(\rho)\ge0$ per active period, let initial administration cost be $A_0$, and let the additional distortionary financing cost of the initial grant be $\lambda s$, with $\lambda\ge0$. Audit resource costs are already in $a$; any extra financing burden for them must either be included there or added explicitly. At the high-quality equilibrium social value is
\[
 W=\frac{v_H-c_H-a(\rho)}{1-\delta}-K-A_0-\lambda s.
 \tag{3}
\]
Fees cancel between clients and the provider. The grant itself cancels between taxpayers and its recipient under equal welfare weights. Its gross fiscal outlay is still $s$, and ongoing audit spending is a separate fiscal requirement. Calling this outcome free of client subsidies does not mean certification has no public operating cost.

For example, take $\delta=0.8$, $\rho=0.5$, $c_H=4$, $c_L=1$, $v_H=7$, and $p=6$. The minimum price in (1) is $5.5$. Provider value is $10$, while the deviation gain is $3$ and the expected continuation penalty is $0.8(0.5)(10)=4$. Set $K=12$, $a(\rho)=0.2$, $A_0=0$, and $\lambda=0.2$. The smallest entry grant is $2$; value (3) is $14-12-0.4=1.6$. Clients continue paying $6$ after the grant ends. All figures are constructed model inputs.

If $K$ instead equals $20$, the minimum grant is $10$ and value becomes $14-20-2=-8$. Provider participation can still be purchased, but this does not establish positive social value. Likewise, a market can have positive $v_H-c_H$ yet no price satisfying (1), for example if $v_H=5$ with the other incentive parameters unchanged.
\section{Finite support cannot manufacture a continuation incentive}
Suppose grants are scheduled only through a deterministic date $T$ and the post-$T$ environment is stationary. At every good history after $T$, the high-quality incentive condition is again (1). If its interval is empty, no strategy promising permanent high quality and stationary prices acceptable to new clients after $T$ can be sustained by the earlier transfers. The reason is local: those transfers are sunk, whereas a low-quality deviation saves $c_H-c_L$ now. This statement does not exclude a subsidy that permanently improves detection, reduces service cost, builds verified reputation, or finances a forfeitable bond; each changes a primitive omitted from the sunk-transfer benchmark.

A useful evaluation would therefore measure the post-support price, repeat demand, verified quality, provider continuation value, and the possibility of reentry after certification loss. Demonstrating more purchases during the supported phase identifies none of these conditions by itself.
\section{Remaining gap and source relationship}
The model fixes client demand, the provider's outside option, audit reliability, the price contract, and a credible exclusion institution. It does not identify those primitives, solve strategic entry with many providers, establish effects on actual business profits, or select among multiple reputation equilibria. Risk aversion and heterogeneous client willingness to pay may alter the feasible price interval. Credit constraints can make (2) insufficient even when participation holds. Welfare weights can make transfers distributively relevant.

The broad catalogue problem remains unresolved. The exact contribution here is the implementability interval and its separation from entry assistance and welfare. The incentive proof is a standard repeated-contract argument supplied in full; no claim of a new general reputation theorem is made.
\begin{thebibliography}{9}
\bibitem{training} D. McKenzie, C. Woodruff, and coauthors (2025).
\emph{Training Entrepreneurs}, VoxDevLit 1(4), Chapter 6.
The primary conclusion identifies functioning training and consulting markets as an evidence gap; it does not supply this repeated-contract model.
\url{https://voxdev.org/voxdevlit/training-entrepreneurs-issue-4/conclusions-training-entrepreneurs}.
\end{thebibliography}

\end{document}
